\documentclass{beamer}
\setbeamerfont{headline}{size=\footnotesize}

\mode<presentation>{
\usetheme{Montpellier}
\setbeamercovered{transparent}
\usecolortheme{lsc}
}

\mode<handout>{
  \usepackage[bar]{beamerthemetree}
  \beamertemplatesolidbackgroundcolor{black!5}
}

\usepackage{amsmath,amssymb}
\usepackage[brazil]{varioref}
\usepackage[english,brazil]{babel}
\usepackage[utf8]{inputenc}
\usepackage{fontawesome}
\usepackage{tikz}
\usetikzlibrary{shapes,arrows,positioning,calc}
\usepackage{graphicx}
\usepackage{listings}
\usepackage{url}
\usepackage{colortbl}
\usepackage{caption}
\usepackage{transparent}
\usepackage{multirow}

\definecolor{pyString}{HTML}{933797}
\definecolor{pyComment}{HTML}{228B22}
\definecolor{pyEmph2}{HTML}{DD52F0}
\definecolor{pyEmph3}{HTML}{D17032}

\lstloadlanguages{Python}
\lstset{
  language=Python,
  basicstyle=\scriptsize\sffamily,
  numberstyle=\color{gray},
  stringstyle=\color{pyString},
  commentstyle=\color{pyComment}\sffamily,
  emph={[2]from,import,pass,return}, emphstyle={[2]\color{pyEmph2}},
  emph={[3]range}, emphstyle={[3]\color{pyEmph3}},
  emph={[4]for,in,def}, emphstyle={[4]\color{blue}},
  showstringspaces=false,
  breaklines=true,
  prebreak=\mbox{{\color{gray}\tiny$\searrow$}},
  numbers=left,
  xleftmargin=15pt
}
%% Custom Commands for Beamer Presentations
%% Este arquivo pode ser usado para definir comandos customizados

% Exemplo de comando customizado (descomente se necessário)
% \newcommand{\alert}[1]{\textcolor{red}{#1}}


\beamertemplatetransparentcovereddynamic
\addtobeamertemplate{title page}{\centering{ \includegraphics[width=1.25cm]{figs/logo.png}}}{}
\title[Estrutura de Dados I]{Recursão}
\author[Elton Sarmanho]{ 
Professor: Elton Sarmanho\inst{1} 
\newline
E-mail: \href{mailto:eltonss@ufpa.br}{eltonss@ufpa.br}\newline
\href{https://github.com/eltonsarmanho}{\faGithub} \href{https://www.linkedin.com/in/elton-sarmanho-836553185/}{\faLinkedinSquare}
}   
\institute[UFPA]{
   \inst{1}%
     Faculdade de Sistemas de Informação - UFPA/CUTINS
}     
 
\date{\today}
\pgfdeclareimage[height= 1.25 cm]{inf}{figs/logo.png}
\logo{{\transparent{0.35}\pgfuseimage{inf}}}

\setbeamertemplate{navigation symbols}{\insertframenavigationsymbol \insertsectionnavigationsymbol \insertbackfindforwardnavigationsymbol} 
\newcounter{saveenumi}
\newcommand{\seti}{\setcounter{saveenumi}{\value{enumi}}}
\newcommand{\conti}{\setcounter{enumi}{\value{saveenumi}}}

\expandafter\def\expandafter\insertshorttitle\expandafter{%
  \insertshorttitle\hfill%
  \insertframenumber\,/\,\inserttotalframenumber}

\begin{document}
\setbeamertemplate{logo}{}
\begin{frame}
\titlepage
\end{frame}
%Começar mostrar Logo apos slide de titulo
\logo{{\transparent{0.35}\pgfuseimage{inf}}}

\begin{frame}
\frametitle{Roteiro}
\tableofcontents[sections={1-2}]
\end{frame}

\section{Conceitos Gerais}
\subsection{Objetivos}
\begin{frame}{Objetivos}
\begin{itemize}
    \item \textbf{Propósito geral:} Apresentar a recursividade e sua potencialidade na resolução de problemas complexos na Ciência da Computação.
    \pause
    \item Ao final desta aula, o aluno será capaz de:
    \begin{itemize}
        \item Definir rigorosamente os conceitos de \textbf{caso base} e \textbf{passo indutivo} em algoritmos.
        \pause
        \item Distinguir cenários ideais para aplicar a estratégia recursiva em problemas práticos.
        \pause
        \item Formular relações de \textbf{recorrência} para delimitar complexidade (Abordagem de Divisão e Conquista).
        \pause
        \item Compreender desafios reais, incluindo recursões implícitas e problemas relativos à pilha de execução (Stack Overflow).
    \end{itemize}
\end{itemize}
\end{frame}

\subsection{Definição}
\begin{frame}{O que é Recursão?}
\begin{block}{Na Matemática}
    Consiste em definir um objeto matemático (função, sequência ou conjunto) em termos da sua própria definição preestabelecida de forma encapsulada.
\end{block}
\pause
\begin{block}{Na Computação}
    É o mecanismo pelo qual uma sub-rotina (função, método procedimental) invoca intencionalmente a \textbf{si mesma} visando solucionar subproblemas progressivamente menores, até satisfazer uma condição trivial preexistente.
\end{block}
\pause
\vspace{0.3cm}

\end{frame}

\begin{frame}{Pilares da Recursividade em Algoritmos}
O avanço para a resolução de sistemas recursivos repousa necessariamente em duas bases axiomáticas:
\pause
\vspace{0.3cm}

\begin{alertblock}{1. Caso Base (Condição de Ancora)}
Demarcação onde o limite da rotina é alcançado. Quando atingida, emite uma resposta computável prontamente (convergência real), assegurando encerramento.
\end{alertblock}
\pause
\begin{block}{2. Passo Indutivo}
Caminho analítico reduzido de instâncias macro. A cada passo, o algoritmo gera dados que tendem iterativamente até encontrar o caso base.
\end{block}
\end{frame}

\begin{frame}[fragile]{Exercício Lógico: Fatorial Matemático ($n!$)}

Para dominar recursividade escalar, extraímos o primórdio produto de cadeias: 
\begin{block}{Definição Recursiva Equacional}
  \begin{equation*}
  n! =
    \begin{cases}
      1, & \text{se } n = 0 \quad \text{(\textbf{Caso Base})}  \\      
      n \cdot (n-1)!, & \text{se } n > 0 \quad \text{(\textbf{Passo de Indução})}
    \end{cases}       
  \end{equation*}
\end{block}
\pause
\textbf{Implementação Computacional:}
\begin{verbatim}
def fatorial(n):
    if n == 0:
        return 1 
    else:
        return n * fatorial(n - 1)  
\end{verbatim}
\end{frame}

\begin{frame}{Tratamento de Contexto: Pilha de Execução ({\it Call Stack})}
Como o sistema preserva variáveis locais entre invocações? Todo compilador provisiona uma \textbf{Pilha (Stack)} para estruturar sucessões não resolvidas.
\vspace{0.2cm}
\begin{center}
\begin{tikzpicture}[>=stealth, font=\small,
    stacknode/.style={draw, rectangle, minimum width=3.2cm, minimum height=0.55cm, fill=blue!10}]
    
    % Base
    \draw[thick, fill=gray!20] (-2.2,0) rectangle (2.2, 0.2);
    
    \node<2->[stacknode] (s1) at (0, 0.7) {\texttt{fatorial(3)} $\rightarrow 3 \times \textbf{?}$};
    \node<3->[stacknode] (s2) at (0, 1.4) {\texttt{fatorial(2)} $\rightarrow 2 \times \textbf{?}$};
    \node<4->[stacknode] (s3) at (0, 2.1) {\texttt{fatorial(1)} $\rightarrow 1 \times \textbf{?}$};
    \node<5->[stacknode, fill=green!20] (s4) at (0, 2.8) {\texttt{fatorial(0)} $\rightarrow 1$};
    
    % Calls
    \draw<3->[->, thick, blue] (2.0, 0.7) .. controls (2.6,0.9) and (2.6,1.2) .. (2.0, 1.4) node[midway, right] {chama};
    \draw<4->[->, thick, blue] (2.0, 1.4) .. controls (2.6,1.6) and (2.6,1.9) .. (2.0, 2.1);
    \draw<5->[->, thick, blue] (2.0, 2.1) .. controls (2.6,2.3) and (2.6,2.6) .. (2.0, 2.8);
    
    % Returns
    \draw<8->[<-, thick, red] (-2.0, 0.7) .. controls (-2.6,0.9) and (-2.6,1.2) .. (-2.0, 1.4) node[midway, left] {rt: $6$};
    \draw<7->[<-, thick, red] (-2.0, 1.4) .. controls (-2.6,1.6) and (-2.6,1.9) .. (-2.0, 2.1) node[midway, left] {rt: $2$};
    \draw<6->[<-, thick, red] (-2.0, 2.1) .. controls (-2.6,2.3) and (-2.6,2.6) .. (-2.0, 2.8) node[midway, left] {rt: $1$};
    
\end{tikzpicture}
\end{center}
\end{frame}

\begin{frame}{Múltipla Ramificação: Sequência de Fibonacci}
Cálculos recursivos podem propagar expansões laterais massivas. Fibonacci exemplifica ramificação de profundidade.

\vspace{0.2cm}
\begin{block}{Relação de Recorrência}
  \begin{equation*}
   f(n) =
    \begin{cases}
      0, & \text{se } n = 0 \quad \text{(\textbf{Base A})}  \\
      1, & \text{se } n = 1 \quad \text{(\textbf{Base B})}  \\
      f(n-1) + f(n-2), & \text{se } n > 1 \quad \text{(\textbf{Bifurcação de Passo})}
    \end{cases}        
  \end{equation*}
\end{block}
\pause
\begin{itemize}
  \item Termos avaliados ($f(0)$ em diante):
  $$0, 1, 1, 2, 3, 5, 8, 13, 21, 34, 55 \dots$$
\end{itemize}
\end{frame}

\begin{frame}{Desdobramento Assintótico: Árvore Fibonacci}
Demonstrativo do Overlap (Superposição) Computacional $\implies O(2^n)$.
\vspace{0.3cm}
\begin{center}
\begin{tikzpicture}[level distance=1.2cm, 
   level 1/.style={sibling distance=3.5cm},
   level 2/.style={sibling distance=1.5cm},
   level 3/.style={sibling distance=0.8cm}, scale=0.9, transform shape]
   
\node {\texttt{f(4)}}
  child {node {\texttt{f(3)}}
    child {node {\texttt{f(2)}}
      child {node {\texttt{f(1)}}}
      child {node {\texttt{f(0)}}}
    }
    child {node[fill=red!20] {\texttt{f(1)}}}
  }
  child {node {\texttt{f(2)}}
    child {node[fill=red!20] {\texttt{f(1)}}}
    child {node[fill=gray!20] {\texttt{f(0)}}}
  };
\end{tikzpicture}
\end{center}

\end{frame}


\subsection{Problemas em Recursividade}
\begin{frame}{Avaliação Qualitativa: Omissões e Otimizações}


\begin{itemize}
 \item \textbf{Recursos de Cauda (\textit{Tail Recursion}):} Se a etapa induzida representar fundamentalmente a encerramento natural avaliado pelo ciclo isolado, sistemas contemporâneos convertem nativamente sob $O(1)$ \textit{em memórias} contornando a pilha.
 \item \textbf{Vícios Axiomáticos:} Abertura recursiva com condicionalidade submetida incorretamente resulta por si em loop infinito.
\end{itemize}
\end{frame}


\subsection{Trabalho 1 - Recursividade}
\begin{frame}[allowframebreaks]{Trabalho Avaliativo - Domínio Modular}
\begin{enumerate}
 \item \textbf{Múltiplo Escalamento Marginal:} Implemente uma diretriz computacional rigorosamente recursiva visando recuperar sem iteração paralela, o mínimo absoluto (\text{menor elemento}) alocado sobre um grande vetor desornado.
 
 \item \textbf{Segmentações Restritas:} Formule uma assinatura capaz de renderizar perfeitamente todas as partições (\textit{substrings contíguas}) provenientes de primitivas sequenciais puras. \\
 Exemplo isolado à constante $S \equiv \text{dog}$:
 $$| d | \ do \ | \ dog \ | \ o \ | \ og \ | \ g |$$
 
 \item \textbf{Manipulação Indutiva Discreta:} Analise a sucessão processual subjacente a representação $g_1 \dots g_n$ perante um parâmetro demarcador universal $k$:
  \begin{equation}
  g_j =
    \begin{cases}
      j-1 & \textbf{se } \: 1 \leq j \leq k   \\
      g_{j-1}+g_{j-2} & \textbf{se } \: j > k  \\
    \end{cases}       
  \end{equation}
 Elabore um algoritmo logico com base nessa expressão matemática.
\end{enumerate}
\end{frame}


\section{Recorrências}
\subsection{Conceito Computacional}
\begin{frame}{Compreendendo Recorrências em Algoritmos}    
\footnotesize
A teoria analítica repousa fortemente no dimensionamento explícito de processos autodescritos via formulações restritivas.

\begin{block}{Formalização Clássica}
Consistem de formulações matemáticas parametrizadas regidas sobre a entrada que especificam \textbf{tempo ou recursos limitantes} expressos nos domínios marginais subsequentes da própria recursão:
$$T(n) = T(n-1) + \Theta(1)$$
\end{block}
\pause
\begin{itemize}
 \item \textbf{Presença Essencial:} Predominam na validação métrica da tríade recursiva (Dividir, Resolver, Combinar) como MergeSort.
 \item \textbf{Abordagens Propostas:} Visam isolar matematicamente expressões absolutas $\Theta(\dots)$ eliminando chamadas cíclicas relativas ao próprio \textit{T()}.
\end{itemize}
\end{frame}


\begin{frame}{Gabaritos Universais: Classificadores Clássicos}    
  \footnotesize

Revisitando matrizes temporais padrão de mercado na Ciência Computacional Básica:
\pause
\begin{description}
 \item[$T(n) = T(n-1) + \Theta(1) \implies O(n)$] 
 \hfill \break Extração serial exaure apenas passo elementar (Buscas desordenadas).
 \pause
 \item[$T(n) = T(n-1) + n \implies O(n^2)$] 
 \hfill \break Varreduras residuais lineares repetidas perante extração paralela fracionada.
 \pause
 \item[$T(n) = T(n/2) + c \implies \Theta(\log n)$] 
 \hfill \break Segmentação simétrica com triagem binária restrita constante (Buscadores Diretos).
 \pause
 \item[$T(n) = 2T(n/2) + n \implies \Theta(n \log n)$] 
 \hfill \break Ramificações paralelas onde as fusões custam um empenho linear absoluto sobre o nível arbóreo correspondente.
\end{description}
\end{frame}

\subsection{Modelagem: Busca Binária}
\footnotesize

% \begin{frame}{Pesquisa Binária (Dicotômica): Transcrição Analítica}    
% Ao procurarmos $x$ uniformizado sobre as instâncias $A[lo \dots hi]$, observamos a restrição radical a uma meta seccionada.

% \vspace{0.4cm}
% \begin{center}
% \begin{tikzpicture}[font=\small, scale=0.8, transform shape]
%     % Array full
%     \draw (0,0) grid (10,1);
%     \node at (5, 1.3) {Amostragem Bruta Dimensional: $n$ ($lo \dots hi$)};
    
%     % Array split
%     \draw<2->[fill=blue!10] (0,-1.5) grid (5,-0.5);
%     \draw<2->[fill=gray!20] (5,-1.5) grid (10,-0.5);
%     \draw<2->[->, thick, red] (5,-0.5) -- (5,0);
%     \node<2-> at (5, -0.2) {$mid = \lfloor \frac{lo+hi}{2} \rfloor$};
    
%     % Narrow down
%     \node<3->[align=left, text width=7cm] at (5, -2.5) {Considerando meta estatística inferior as margens preestabelecidas sobre o $mid$, \textbf{descartamos e focamos restritamente lateral esquerda} da árvore de dados reduzida perfeitamente pela aproximação fator $\frac{n}{2}$.};
    
%     \draw<4->[fill=red!20] (2.5,-4) grid (5,-3);
%     \draw<4->[->, thick] (2.5,-3) -- ++(0,1);
% \end{tikzpicture}
% \end{center}

% \vspace{0.3cm}
% \onslide<5->{
% \textbf{Modelo Gerado:}  $T(n) = T(n/2) + c \quad \text{e} \quad T(1) = 1$
% }
% \end{frame}


\subsection{Método da Iteração}
\begin{frame}{Método Avançado por Substituição (Desdobramento)}    
\footnotesize

Baseia-se no colapso paramétrico gradual almejando identificar séries progressivas subjacentes. Assumindo modelo primário: $\mathbf{T(n) = T\left(\frac{n}{2}\right) + c}$

\pause
\vspace{0.3cm}
\textbf{Expansão Passo-a-Passo:}
\begin{align*}
 T(n) &= \mathbf{T(n/2)} + c \onslide<3->{= \mathbf{\left[T(n/4) + c\right]} + c} \\
      &\onslide<4->{= \mathbf{T(n/2^2) + 2c}} \\
      &\onslide<5->{= \mathbf{\left[T(n/8) + c\right]} + 2c = \mathbf{T(n/2^3) + 3c}} \\
      &\onslide<6->{\implies \dots \\}
      &\onslide<7->{= \mathbf{T(n/2^k) + c \cdot k} \quad \text{(Equação Genérica de pulso } k\text{)}}
\end{align*}
\end{frame}


\begin{frame}{Isolamento da Variável $k$ em Escala Base}    
\footnotesize

Assumindo que nossa âncora de saída (Condição Base) processa arranjos iteracionais triviais de tamanho 1, inferimos $\mathbf{T(1) = 1}$.
\pause
\begin{block}{Buscando convergência $n/2^k = 1$}
\begin{align*}
  \frac{n}{2^k} &= 1 \onslide<3->{\ \implies\  n = 2^k} \\
  &\onslide<4->{\implies\  \log_2 (n) = \log_2 (2^k)} \\
  &\onslide<5->{\implies\  k = \log_2 n}
\end{align*}
\end{block}

\pause[6]
Retornando o valor $k$, finalizamos para as iterações da Pesquisa:
$$T(n) = T(1) + c \cdot (\log_2 n)$$
\begin{alertblock}{Conclusão Assintótica}
O Algoritmo opera quantitativamente por custo marginal rigoroso $\mathbf{\Theta(\log n)}$.
\end{alertblock}
\end{frame}


\subsection{Árvore de Recursão}
\begin{frame}{Projeção Fractal: Árvore de Recursão Clássica}    

Ferramenta visualmente proeminente aplicável aos \textit{design} de algoritmos Divisão-e-Conquista modelados na essência $\mathbf{T(n) = 2T(n/2) + \Theta(n)}$.
\pause
\begin{center}
\begin{tikzpicture}[level distance=1.1cm, 
   level 1/.style={sibling distance=3.5cm},
   level 2/.style={sibling distance=1.5cm},
   level 3/.style={sibling distance=1cm}, scale=0.85, transform shape,
   every node/.style={circle, minimum size=0.6cm}]
   
\node {$n$}
  child {node {$\frac{n}{2}$}
    child {node {$\frac{n}{4}$}
      child[dotted] {node[draw=none]{}}
      child[dotted] {node[draw=none]{}}
    }
    child {node {$\frac{n}{4}$}
      child[dotted] {node[draw=none]{}}
      child[dotted] {node[draw=none]{}}
    }
  }
  child {node {$\frac{n}{2}$}
    child {node {$\frac{n}{4}$}
      child[dotted] {node[draw=none]{}}
      child[dotted] {node[draw=none]{}}
    }
    child {node {$\frac{n}{4}$}
      child[dotted] {node[draw=none]{}}
      child[dotted] {node[draw=none]{}}
    }
  };

% Right labels
\node<2->[draw=none, right=3cm] at (0, 0) {Custo: $n$};
\node<3->[draw=none, right=1.25cm] at (1.75, -1.1) {Custo: $\frac{n}{2} + \frac{n}{2} = n$};
\node<4->[draw=none, right=0cm] at (2.6, -2.2) {Custo: $4 \times \frac{n}{4} = n$};
\node<5->[draw=none, right=0.3cm] at (2.5, -3.3) {Nível Secundário...};

\end{tikzpicture}
\end{center}
\onslide<6->{Somatório cumulativo $\sum^{h}_{1} {n} = n \times h$ onde a profundidade foliar $h = \log_2 n$.
\textbf{Solução de Complexidade:} $\mathbf{\Theta(n \log n)}$.}

\end{frame}


\subsection{Método Master}
 \begin{frame}{Ensaio de Master: Formulário Direto Mapeado}   

Trata-se do \textit{Cookbook} algorítmico global \cite{A03} visando solucionar eficientemente recorrências regidas pelo esqueleto genérico base:
\begin{alertblock}{Estrutura do Teorema}
 $$\mathbf{T(n) = aT\left(\frac{n}{b}\right) + f(n)} \quad \text{em que } a\geq 1, b > 1 \text{ ctes, e } f(n) \geq 0$$
\end{alertblock}
\pause
\textbf{Decifrando Parâmetros da Lógica Analítica:}
\begin{itemize}
 \item $a$: Subproblemas fracionados disparados ativamente recursão paralela na máquina.
 \item $b$: Fator de cisão/fatiamento reduzindo o problema nuclear.
 \item $f(n)$: Fator de carga de custo para repartição central e fusão/mesclagem subsequente do processo.
\end{itemize}
\end{frame}


\begin{frame}{Avaliação de Domínio: Comparativo Polinomial}   
O pilar metodológico se encarrega de comparar assintoticamente o esforço polinomial fixo $f(n)$ contra a expansão quantizada $n^{\log_b a}$ (volume iterativo da árvore tracionada).

\pause
\begin{description}
 \item[Caso 1: Folhas predominam] Se $f(n) = O(n^{\log_b a - \epsilon})$ com $\epsilon > 0 \implies \mathbf{T(n) = \Theta(n^{\log_b a})}$. \\ \textit{(A recursão cresce velozmente, devorando recursos marginais na borda foliar).}
 \pause
 \item[Caso 2: Empate de Forças] Se $f(n) = \Theta(n^{\log_b a}) \implies \mathbf{T(n) = \Theta(n^{\log_b a} \cdot \log n)}$. \\ \textit{(Trabalho equitativo sendo distribuído ao longo dos níveis perfeitamente).}
 \pause
 \item[Caso 3: Raiz predominante] Se $f(n) = \Omega(n^{\log_b a + \epsilon})$ para $\epsilon > 0$, dado a regularidade restrita $a \cdot f(\frac{n}{b}) \leq c \cdot f(n) \implies \mathbf{T(n) = \Theta(f(n))}$.
\end{description}
\end{frame}


\begin{frame}{Casos de Aplicação Direta Mestre}   
Embate Intelectual Clássico - Tente determinar o limite Assintótico:
\pause
\begin{block}{Exemplo 1 (Raízes Duplas)}
$T(n) = 4T(n/2) + n \implies n^{\log_2 4} = n^2$. Folhas sobrepujam $n$. \\
Resposta: \textbf{Caso 1 $\rightarrow \Theta(n^2)$}.
\end{block}
\pause
\begin{block}{Exemplo 2 (Esboço Singular)}
$T(n) = T(2n/3) + 1 \implies b=3/2 \implies n^{\log_{1.5} 1} = n^0 = 1$. Concorrem equilibrados \\
Resposta: \textbf{Caso 2 $\rightarrow \Theta(\log n)$}.
\end{block}
\pause
\begin{block}{Exemplo 3 (Carga Complexificada)}
$T(n) = 9T(n/3) + n^3 \implies n^{\log_3 9} = n^2$. Raiz (composição) supera limite $n^3 > n^2$ \\
Resposta: \textbf{Caso 3 $\rightarrow \Theta(n^3)$}.
\end{block}
\end{frame}


\section{Referências Bibliográficas}
\begin{frame}[t,allowframebreaks]
\frametitle{Referências}
\begin{thebibliography}{*}

        \bibitem[A01]{A01} Lee K.D., Hubbard S. (2015) Trees. In: Data Structures and Algorithms with Python. Undergraduate Topics in Computer Science. Springer, Cham. \url{https://doi.org/10.1007/978-3-319-13072-9_6}
        
        \bibitem[A02]{A02} Hubbard, J. (2007). Schaum’s Outline of Data Structures with Java. \url{http://www.amazon.com/Schaums-Outline-Data-Structures-Java/dp/0071476989}
        
        \bibitem[A03]{A03} Cormen, T. H., Leiserson, C. E., \& Stein, R. L. R. (2012). Algoritmos: teoria e prática. \url{https://books.google.com.br/books?id=6iA4LgEACAAJ}.
        
        \bibitem[A04]{A04} Ascenio, Ana Fernanda Gomes. Estrutura de dados: Algoritmos, análise da complexidade e implementações em Java e C++. São Paulo: Pearson Prentice Hall, 2010.
        
        \bibitem[A06]{A06} Szwarcfiter, Jayme Luiz. Estruturas de dados e seus algoritmos / Jayme Luiz Szwarcfiter, Lilian Markenzon. 3.ed. [Reimpr.]. - Rio de Janeiro : LTC, 2015.
        
        \bibitem[A07]{A07} Capítulo 20: Árvores — documentação Aprenda Computação com Python 3.0 2009.1. \url{https://mange.ifrn.edu.br/python/aprenda-com-py3/capitulo_20.html}
\end{thebibliography}
\end{frame}

\section{}
\setbeamertemplate{logo}{}
\begin{frame}
\titlepage
\end{frame}

\end{document}
